\documentclass[a4paper,11pt]{letter}
% POZOR: samo utf8. V starih datotekah je bil za njim še \usepackage[cp1250]{inputenc},
% kar je prevzelo prednost in lahko pokvari šumnike.
\usepackage[utf8]{inputenc}
\usepackage[T1]{fontenc}
\usepackage{lmodern}
\usepackage[english,slovene]{babel}
\usepackage[tikz]{bclogo}
\usepackage{qrcode,marvosym}
\usepackage{pgf,tikz,pgfplots}
\pgfplotsset{compat=1.14}
\usepackage{mathrsfs}
\usetikzlibrary{arrows}
\usepackage{graphicx}
\newcommand{\ds}{\displaystyle}
\usepackage{fancybox}
\usepackage{amsfonts}
\usepackage{amsmath}
\usepackage{wallpaper}
\usepackage{hyperref}
\usepackage{array}
\newcolumntype{M}[1]{>{\centering\arraybackslash}m{#1}}
\newcolumntype{N}{@{}m{0pt}@{}}
\usepackage{fancyhdr,enumerate}
\newcommand{\ora}{\overrightarrow}
\renewcommand{\headrulewidth}{3pt}
\renewcommand{\headrule}{\hbox to\headwidth{%
  \color{gray}\leaders\hrule height \headrulewidth\hfill}}
\renewcommand{\footrulewidth}{2pt}
\renewcommand{\footrule}{\hbox to\headwidth{%
  \color{brown}\leaders\hrule height \footrulewidth\hfill}}
\usepackage{gensymb}
\usepackage{amssymb}
  \textwidth 20 true cm
  \oddsidemargin -2 true cm
  \evensidemargin -2 true cm
  \topmargin -3.5 true cm
  \textheight 27.5 true cm
  \linespread{1.2}
\renewcommand{\headrulewidth}{0.3pt}
\renewcommand{\footrulewidth}{1.9pt}
\definecolor{qqwuqq}{rgb}{1,0,0}
\definecolor{qqqqff}{rgb}{0,0,1}
\definecolor{qqffqq}{rgb}{0,1,0}
\definecolor{ffqqqq}{rgb}{1,0,0}
\definecolor{zzuxux}{rgb}{0,0,0}
\usepackage{mdframed}
\usepackage{multicol}
\definecolor{bgblue}{RGB}{0,0,253}
\usepackage{eurosym}
%%%%%%%%%%%%%%%%%%%%%%%%%%%%%%%%%%%%%%%%%%%%%%%%%%%%%%%%
\newcounter{tocke}
\setcounter{tocke}{0}
\usepackage{tikz-3dplot}
%%%%%%%%%%%%%%%%%%%%%%%%%%%%%%%%%%%%%%%%%%%%%%%%%%%%%%%%%%
\mdfdefinestyle{theoremstyle}{%
linecolor=brown!40,linewidth=1pt,%
frametitlerule=true,%
frametitlebackgroundcolor=brown!50,
innertopmargin=\topskip,
}
\mdfdefinestyle{theoremstyle1}{%
linecolor=brown,middlelinewidth=1pt,%
frametitlerule=true,%
apptotikzsetting={\tikzset{mdfframetitlebackground/.append style={%
shade,left color=brown!40, right color=gray!10},
mdfbackground/.append style={
top color=white!100!white,
bottom color=black!2
},
}},
frametitlerulecolor=gray,
frametitlerulewidth=1pt,
innertopmargin=\topskip,
innerbottommargin=\topskip,
}
\mdtheorem[style=theoremstyle1]{naloga}{Naloga}
\mdtheorem[style=theoremstyle]{resitev}{Rešitev}
%%%%%%%%%%%%%%%%%%%%%%%%%%%%%%%%%%%%%%%%%%%%%%%%%%%%%%%%%%
\usepackage[table]{xcolor}
\usepackage[printwatermark]{xwatermark}
\newwatermark[allpages,color=brown!2,angle=45,scale=5,
xpos=-5,ypos=0]{matej.info}
%%%%%%%%%%%%%%%%%%%%%%%%%%%%%%%%%%%%%%%%%%%%%%%%%%%%%%%%%%

% Razred "letter" ne pozna ukaza \section, zato si naredimo svoj naslov.
\newcommand{\razdelek}[1]{%
  \vspace{0.5cm}
  \begin{center}
  \colorbox{brown!25}{\parbox{0.97\textwidth}{\centering\vspace{2pt}\large\textbf{#1}\vspace{2pt}}}
  \end{center}
  \vspace{0.2cm}
}

\begin{document}
\pagestyle{fancy}
\renewcommand\headrulewidth{0pt}
\lhead{}\chead{}\rhead{}
\rfoot{\vspace*{-2\baselineskip}}
\renewcommand{\vec}{\overrightarrow}
\everymath{\displaystyle}

\begin{minipage}{20cm}
\begin{bclogo}[couleur=brown!40,
  arrondi=0,ombre=true, couleurOmbre=gray!50,
logo =\bcplume,
  noborder=true]
{\textrm{{\Large{NALOGE 1.N.1} \Huge{}:
\large{Stereometrija -- telesa} \hfill $G-4$}}}
\end{bclogo}
\end{minipage}
\begin{minipage}{20cm}
\begin{bclogo}[couleur=brown!40,ombre=true, couleurOmbre=gray!50,
  arrondi=0,
logo =\bcplume,
  noborder=true]
{\textsc{{\large{Ime in priimek:}
\vspace{0.2cm} \rule{15cm}{0.2mm}\hfill }}}
\end{bclogo}
\end{minipage}

\vspace{0.3cm}
\small
Vse oznake za telesa so standardne. Če enota ni zapisana, so dolžine v $m$, ploščine v $m^2$
in prostornine v $m^3$. Rezultate zapiši z ustreznimi enotami. Kote računaj do kotne minute natančno.
\normalsize

\razdelek{I. \ Vaje}

%%%%%%%%%%%%%%%%%%%%%%%%%%%%%%%%%%%%%%%%%%%%%%%%%
\begin{naloga}
V valju s površino $56\pi$ m$^2$ meri polmer 4 m.

a) Izračunaj višino valja.

b) Na eno osnovno ploskev valja nalepimo polkroglo, na drugo pa enakostranični stožec.
Izračunaj površino in prostornino sestavljenega telesa.
\end{naloga}

%%%%%%%%%%%%%%%%%%%%%%%%%%%%%%%%%%%%%%%%%%%%%%%%%
\begin{naloga}
Valj:

a) $V=750\pi,\ v=30$; \quad izračunaj $r,P$.

b) $P=432\pi,\ v=6$; \quad izračunaj $r,V$.

c) $S_{pl}=8,\ P=20{,}64$; \quad izračunaj $r,V$.

d) $S_{pl}=4,\ V=12$; \quad izračunaj $r,v$.

e) $2r=v,\ P=192$; \quad izračunaj $r,V$.
\end{naloga}

%%%%%%%%%%%%%%%%%%%%%%%%%%%%%%%%%%%%%%%%%%%%%%%%%
\begin{naloga}
Stožec ($\omega$ je kot med osnovno ploskvijo in stranico stožca):

a) $P=196,\ v=11{,}4$; \quad izračunaj $r,s,V$.

b) $V=623,\ r=7{,}5$; \quad izračunaj $v,s,P$.

c) $P=96\pi,\ r=6$; \quad izračunaj $v,s,V$.

d) $S_{pl}=20\pi,\ S_0=12\pi$; \quad izračunaj $r,V,\omega$.

e) $s=6,\ r=2\sqrt{5}$; \quad izračunaj $v,\omega$.

f) $V=29160\pi,\ r=54$; \quad izračunaj $v,s,P$.

g) $\omega=30^{\circ},\ v=6$; \quad izračunaj $S_{pl}$.
\end{naloga}

%%%%%%%%%%%%%%%%%%%%%%%%%%%%%%%%%%%%%%%%%%%%%%%%%
\begin{naloga}
Enakostranični trikotnik s stranico $a=4\sqrt{3}$ dm zavrtimo okoli stranice za polni kot.
Izračunaj površino in prostornino tako dobljenega telesa.
\end{naloga}

%%%%%%%%%%%%%%%%%%%%%%%%%%%%%%%%%%%%%%%%%%%%%%%%%
\begin{naloga}
a) Kvader ima površino 248 cm$^2$, razmerje robov pa je $2:3:5$.
Izračunaj prostornino in maso, če je gostota telesa 9 kg/dm$^3$.

b) Kvader ima prostornino 364 m$^3$, razmerje robov pa je $2:14:13$.
Izračunaj dolžino telesne diagonale kvadra in ploščino trikotnika, ki ga določajo
ploskovne diagonale kvadra.
\end{naloga}

%%%%%%%%%%%%%%%%%%%%%%%%%%%%%%%%%%%%%%%%%%%%%%%%%
\begin{naloga}
$n$-strana prizma:

a) $n=3,\ a=4,\ b=6,\ c=7,\ v=8$; \quad izračunaj $P,V$.

b) $n=3,\ V=12992,\ a=35,\ b=54,\ c=47$; \quad izračunaj $v,P$.

c) $n=6,\ a=4\sqrt{3}$, pravilna enakoroba; \quad izračunaj $P,V$.

d) osnovna ploskev je paralelogram $a=8,\ b=6,\ \alpha=30^{\circ}$,
višina prizme je enaka krajši diagonali osnovne ploskve; \quad izračunaj $P,V$.

e) osnovna ploskev je romb $e=8,\ f=6$, prizma je enakoroba; \quad izračunaj $P,V$.
\end{naloga}

%%%%%%%%%%%%%%%%%%%%%%%%%%%%%%%%%%%%%%%%%%%%%%%%%
\begin{naloga}
$n$-strana pokončna piramida ($\omega$ je kot med osnovno in stransko ploskvijo,
$\alpha$ je kot med osnovnim in stranskim robom):

a) $n=3,\ a=10,\ V=3608$, pravilna; \quad izračunaj $v,P$ in stranski rob.

b) $n=3,\ P=400\sqrt{3}$, pravilna enakoroba; \quad izračunaj $a,v,V$.

c) $n=3,\ v=8\sqrt{2}$, pravilna enakoroba; \quad izračunaj $P,V$.

d) $n=4,\ a=30$, stranski rob meri 70, pravilna; \quad izračunaj $P,V,\omega,\alpha$.

e) $n=4,\ a=14,\ v=24$, pravilna; \quad izračunaj $P,V,\omega,\alpha$.

f) $n=12,\ a=1$, pravilna, višina je enaka polmeru očrtanega kroga osnovni ploskvi;
\quad izračunaj $P,V,\omega,\alpha$.
\end{naloga}

%%%%%%%%%%%%%%%%%%%%%%%%%%%%%%%%%%%%%%%%%%%%%%%%%
\begin{naloga}
Prostornina pokončne štiristrane piramide meri 4480 cm$^3$. V osnovni ploskvi -- pravokotniku --
je diagonala za 4 cm daljša od prvega osnovnega roba, drugi rob pa meri 24 cm.
Izračunaj višino in površino piramide.
\end{naloga}
\newpage
%%%%%%%%%%%%%%%%%%%%%%%%%%%%%%%%%%%%%%%%%%%%%%%%%
\begin{naloga}
a) Kroglo včrtamo v enakostranični stožec s stranico $5\sqrt{3}$.
Izračunaj površino in prostornino krogle.

b) Kroglo očrtamo kocki z robom $10\sqrt{2}$. Izračunaj razmerje prostornin obeh teles.

c) Kvadru z robovi 3 m, 4 m in 12 m očrtamo kroglo.
Kolikokrat je prostornina krogle večja od prostornine kvadra?
\end{naloga}

%%%%%%%%%%%%%%%%%%%%%%%%%%%%%%%%%%%%%%%%%%%%%%%%%
\begin{naloga}
Pravokotni trikotnik s katetama 12 cm in 5 cm zavrtimo okoli:

a) daljše katete, \quad b) krajše katete, \quad c) hipotenuze.

Katero telo ima največjo prostornino in katero največjo površino?
\end{naloga}

%%%%%%%%%%%%%%%%%%%%%%%%%%%%%%%%%%%%%%%%%%%%%%%%%
\begin{naloga}
Enakokraki trapez z osnovnicama 10 dm in 3 dm ter krakoma 5 dm zavrtimo okoli krajše osnovnice.

a) Koliko znašata prostornina in površina rotacijskega telesa?

b) Bi se prostornina kaj spremenila, če bi vrteli okoli daljše osnovnice? Kaj pa površina?
\end{naloga}

\razdelek{II. \ Naloge s testov 1.0--1.3}

%%%%%%%%%%%%%%%%%%%%%%%%%%%%%%%%%%%%%%%%%%%%%%%%%
\begin{naloga}
V stožcu s stranico 5 m meri površina $24\pi$ m$^2$.

a) Izračunaj polmer, višino in prostornino stožca.

b) Koliko meri kot med stranico in osnovno ploskvijo?

c) Iz stožca izrežemo največjo možno kroglo. Izračunaj polmer te krogle.
\end{naloga}

%%%%%%%%%%%%%%%%%%%%%%%%%%%%%%%%%%%%%%%%%%%%%%%%%
\begin{naloga}
V tristrani prizmi merijo osnovni robovi 8 cm, 15 cm in 17 cm.
Višina prizme je enaka polmeru očrtanega kroga osnovne ploskve.
Izračunaj prostornino in površino prizme.
\end{naloga}

%%%%%%%%%%%%%%%%%%%%%%%%%%%%%%%%%%%%%%%%%%%%%%%%%
\begin{naloga}
V pravilni štiristrani piramidi meri diagonala osnovne ploskve $4\sqrt{2}$ cm,
vsota vseh stranskih robov pa je $24\sqrt{2}$ cm.

a) Izračunaj rob osnovne ploskve in višino.

b) Izračunaj kot med stranskim in osnovnim robom.

c) Prisekano piramido dobimo tako, da piramido presekamo z ravnino, vzporedno osnovni ploskvi,
1 cm pod vrhom. Koliko meri prostornina prisekane piramide?
\end{naloga}

%%%%%%%%%%%%%%%%%%%%%%%%%%%%%%%%%%%%%%%%%%%%%%%%%
\begin{naloga}
Polkrog s polmerom 5 dm zavrtimo okoli premera za $270^{\circ}$.

a) Izračunaj površino tako dobljene vrtenine.

b) Kolikšna je masa vrtenine, če je gostota snovi $3$ kg/dm$^3$?

c) Telo pretopimo v kocko. Izračunaj dolžino telesne diagonale kocke.
\end{naloga}

%%%%%%%%%%%%%%%%%%%%%%%%%%%%%%%%%%%%%%%%%%%%%%%%%
\begin{naloga}
Štiristrana pokončna prizma ima za osnovno ploskev paralelogram, v katerem sta dolžini stranic
$a=33$ cm in $b=41$ cm, diagonala pa meri 58 cm. Plašč prizme meri 1480 cm$^2$.

a) Izračunaj višino prizme.

b) Koliko meri prostornina prizme?

c) Izračunaj kot med najdaljšo telesno diagonalo in osnovno ploskvijo.
\end{naloga}

%%%%%%%%%%%%%%%%%%%%%%%%%%%%%%%%%%%%%%%%%%%%%%%%%
\begin{naloga}
Enakokraki trapez z osnovnicama 10 dm in 4 dm ter krakoma 5 dm zavrtimo okoli krajše osnovnice.

a) Koliko znašata prostornina in površina rotacijskega telesa?

b) Izračunaj polmer polkrogle, ki ima enako površino kot vrtenina.
\end{naloga}

%%%%%%%%%%%%%%%%%%%%%%%%%%%%%%%%%%%%%%%%%%%%%%%%%
\begin{naloga}
V pravilni štiristrani pokončni piramidi meri rob osnovne ploskve 14 cm, višina pa 24 cm.

a) Izračunaj površino in prostornino telesa.

b) Izračunaj kot med osnovnim in stranskim robom.

c) Izračunaj kot med osnovno in stransko ploskvijo.

d) Prisekano piramido dobimo tako, da piramido presekamo z ravnino, vzporedno osnovni ploskvi,
12 cm pod vrhom. V kakšnem razmerju sta prostornini prisekane piramide in dopolnilne piramide?
\end{naloga}

%%%%%%%%%%%%%%%%%%%%%%%%%%%%%%%%%%%%%%%%%%%%%%%%%
\begin{naloga}
a) Kvadru s stranicami 8 cm, 6 cm in 24 cm očrtamo kroglo. Izračunaj polmer krogle.

b) Iz največje ploskve mreže tega kvadra izrežemo plašč valja in ga zvijemo po daljši stranici.
Koliko meri polmer valja?
\end{naloga}

%%%%%%%%%%%%%%%%%%%%%%%%%%%%%%%%%%%%%%%%%%%%%%%%%
\begin{naloga}
Prečni presek stožca je enakokraki trikotnik s ploščino 12 dm$^2$ in stranico 5 dm.

a) Izračunaj površino in prostornino stožca. \emph{(Nalogo imata dve rešitvi.)}

b) Koliko meri središčni kot v ravnino iztegnjenega plašča stožca?
\end{naloga}

%%%%%%%%%%%%%%%%%%%%%%%%%%%%%%%%%%%%%%%%%%%%%%%%%
\begin{naloga}
V enakostranični valj s prostornino $250\pi$ cm$^3$ včrtamo kroglo.

a) Izračunaj prostornino krogle.

b) Koliko odstotkov površine valja znaša površina krogle?
\end{naloga}

%%%%%%%%%%%%%%%%%%%%%%%%%%%%%%%%%%%%%%%%%%%%%%%%%
\begin{naloga}
Izračunaj prostornino in površino pravilne 8-strane prizme z osnovnim robom $a=6$ dm in višino,
ki je enaka polmeru očrtanega kroga osnovne ploskve. Rezultat zaokroži na tri decimalna mesta.
\end{naloga}

%%%%%%%%%%%%%%%%%%%%%%%%%%%%%%%%%%%%%%%%%%%%%%%%%
\begin{naloga}
V pravilni štiristrani piramidi je osnovna ploskev pravokotnik s stranicama 24 cm in 10 cm,
prostornina pa meri 400 cm$^3$.

a) Izračunaj višino piramide.

b) Koliko meri plašč piramide? Zapiši točen rezultat.

c) Izračunaj stranski rob piramide in kot med osnovno ploskvijo in stranskim robom.
\end{naloga}

%%%%%%%%%%%%%%%%%%%%%%%%%%%%%%%%%%%%%%%%%%%%%%%%%
\begin{naloga}
Romb z diagonalama 24 cm in 10 cm zavrtimo okoli stranice za polni kot.
Izračunaj površino in prostornino vrtenine.
\end{naloga}


%%%%%%%%%%%%%%%%%%%%%%%%%%%%%%%%%%%%%%%%%%%%%%%%%
\begin{naloga}
V tristrani prizmi merijo robovi osnovne ploskve 9 m, 10 m in 17 m.

a) Koliko odstotkov prostornine prizme zavzame valj, ki je včrtan prizmi?

b) Izračunaj površino prizme, ki ima za višino polmer včrtanega kroga osnovni ploskvi.
\end{naloga}
\newpage

%%%%%%%%%%%%%%%%%%%%%%%%%%%%%%%%%%%%%%%%%%%%%%%%%
\begin{naloga}
a) Izračunaj prostornino pravilne pokončne enakorobe štiristrane piramide,
če meri dolžina roba piramide $a=3\sqrt{2}$ cm.

b) Piramido presekamo na polovici višine. Izračunaj površino spodnjega dela telesa.
\end{naloga}

%%%%%%%%%%%%%%%%%%%%%%%%%%%%%%%%%%%%%%%%%%%%%%%%%
\begin{naloga}
Polkrogla in stožec imata skupno osnovno ploskev in enako prostornino.

a) V kakšnem razmerju sta polmer stožca in stranica stožca?

b) V kakšnem razmerju sta površini polsfere in plašča stožca?
\end{naloga}

%%%%%%%%%%%%%%%%%%%%%%%%%%%%%%%%%%%%%%%%%%%%%%%%%
\begin{naloga}
V posodo oblike kvadra, ki ima osnovno ploskev z dimenzijama $8\pi$ dm $\times$ $4\pi$ dm,
potopimo valj s polmerom 12 cm in višino $12\pi$ cm.

a) Za koliko se dvigne voda v posodi?

b) Izračunaj največjo dolžino tanke palice, ki jo še lahko v celoti potopimo v to posodo,
v kateri je višina tekočine $\pi$ dm.
\end{naloga}

\vfill
\begin{center}
\qrcode[hyperlink,height=0.8in]{http://www.matej.info}
\end{center}

\newpage
\razdelek{Končne rešitve}

\small
\begin{enumerate}
\item a.) $v=3$ m \quad
      b.) $P=88\pi\approx 276{,}5$ m$^2$; \ $V=\dfrac{16\pi\left(17+4\sqrt{3}\right)}{3}\approx 400{,}9$ m$^3$

\item a.) $r=5$, $P=350\pi\approx 1099{,}6$ \quad
      b.) $r=12$, $V=864\pi\approx 2714{,}3$ \quad
      c.) $r\approx 1{,}418$, $V\approx 5{,}67$ \quad
      d.) $r=6$, $v=\dfrac{1}{3\pi}\approx 0{,}106$ \quad
      e.) $r=\sqrt{\dfrac{32}{\pi}}\approx 3{,}192$, $V\approx 204{,}26$

\item a.) $r\approx 3{,}909$, $s\approx 12{,}052$, $V\approx 182{,}41$ \quad
      b.) $v\approx 10{,}576$, $s\approx 12{,}966$, $P\approx 482{,}21$ \quad
      c.) $s=10$, $v=8$, $V=96\pi\approx 301{,}59$ \quad
      d.) $r=2\sqrt{3}$, $V=\dfrac{32\sqrt{3}\,\pi}{3}\approx 58{,}04$, $\omega=53^{\circ}7'48''$ \quad
      e.) $v=4$, $\omega=41^{\circ}48'37''$ \quad
      f.) $v=30$, $s=6\sqrt{106}\approx 61{,}77$, $P=\left(2916+324\sqrt{106}\right)\pi\approx 19640{,}6$ \quad
      g.) $r=6\sqrt{3}$, $s=12$, $S_{pl}=72\sqrt{3}\,\pi\approx 391{,}78$

\item Telo je sestavljeno iz dveh skladnih stožcev s skupno osnovno ploskvijo,
      $r=6$ dm, $s=4\sqrt{3}$ dm:
      $P=48\sqrt{3}\,\pi\approx 261{,}19$ dm$^2$, \ $V=48\sqrt{3}\,\pi\approx 261{,}19$ dm$^3$

\item a.) robovi $4$, $6$ in $10$ cm; $V=240$ cm$^3$, $m=2160$ g $=2{,}16$ kg \quad
      b.) robovi $2$, $14$ in $13$ m; $d=\sqrt{369}=3\sqrt{41}\approx 19{,}21$ m,
      ploščina trikotnika $=\dfrac{1}{2}\sqrt{a^2b^2+b^2c^2+c^2a^2}\approx 92{,}98$ m$^2$

\item a.) $S_0\approx 11{,}977$, $P\approx 159{,}95$, $V\approx 95{,}81$ \quad
      b.) $S_0\approx 812{,}24$, $v\approx 16$, $P\approx 3799{,}8$ \quad
      c.) $S_0=72\sqrt{3}$, $P=288+144\sqrt{3}\approx 537{,}42$, $V=864$ \quad
      d.) krajša diagonala $=\sqrt{100-48\sqrt{3}}\approx 4{,}106$, $P\approx 162{,}98$, $V\approx 98{,}55$ \quad
      e.) $a=5$, $P=148$, $V=120$

\item a.) $v=250$, $P\approx 3793{,}6$, stranski rob $\approx 250{,}07$ \quad
      b.) telo je pravilni tetraeder: $a=20$, $v=\dfrac{20\sqrt{6}}{3}\approx 16{,}33$,
      $V=\dfrac{2000\sqrt{2}}{3}\approx 942{,}81$ \quad
      c.) $a=8\sqrt{3}$, $P=192\sqrt{3}\approx 332{,}55$, $V=128\sqrt{6}\approx 313{,}53$ \quad
      d.) $v=\sqrt{4450}\approx 66{,}708$, $P\approx 5002{,}4$, $V\approx 20012{,}5$,
      $\omega=77^{\circ}19'38''$, $\alpha=72^{\circ}21'34''$ \quad
      e.) $P=896$, $V=1568$, $\omega=73^{\circ}44'23''$, $\alpha=67^{\circ}35'6''$ \quad
      f.) $P\approx 27{,}31$, $V\approx 7{,}21$, $\omega=45^{\circ}59'35''$, $\alpha=45^{\circ}$

\item prvi rob meri 70 cm, $v=8$ cm, $P\approx 3551{,}2$ cm$^2$

\item a.) $\varrho=\dfrac{5}{2}$, $P=25\pi\approx 78{,}54$, $V=\dfrac{125\pi}{6}\approx 65{,}45$ \quad
      b.) $R=5\sqrt{6}$, \ $V_{krogle}:V_{kocke}=\dfrac{\sqrt{3}\,\pi}{2}\approx 2{,}72$ \quad
      c.) $R=\dfrac{13}{2}$, \ $V_{krogle}:V_{kvadra}=\dfrac{2197\pi}{864}\approx 7{,}99$

\item a.) $V=100\pi\approx 314{,}16$ cm$^3$, $P=90\pi\approx 282{,}74$ cm$^2$ \quad
      b.) $V=240\pi\approx 753{,}98$ cm$^3$, $P=300\pi\approx 942{,}48$ cm$^2$ \quad
      c.) $r=\dfrac{60}{13}$, $V=\dfrac{1200\pi}{13}\approx 289{,}99$ cm$^3$,
      $P=\dfrac{1020\pi}{13}\approx 246{,}49$ cm$^2$ \\
      Največjo prostornino \emph{in} največjo površino ima telo iz primera b.).

\item a.) $v=\dfrac{\sqrt{51}}{2}\approx 3{,}571$ dm, $V=\dfrac{391\pi}{4}\approx 307{,}09$ dm$^3$,
      $P=15\sqrt{51}\,\pi\approx 336{,}53$ dm$^2$ \quad
      b.) Da, oboje se spremeni: $V=68\pi\approx 213{,}63$ dm$^3$,
      $P=8\sqrt{51}\,\pi\approx 179{,}48$ dm$^2$

\item a.) $r=3$ m, $v=4$ m, $V=12\pi\approx 37{,}70$ m$^3$ \quad
      b.) $53^{\circ}7'48''$ \quad
      c.) $\varrho=1{,}5$ m

\item Osnovna ploskev je pravokotni trikotnik ($8^2+15^2=17^2$): $S_0=60$ cm$^2$, $v=8{,}5$ cm,
      $V=510$ cm$^3$, $P=460$ cm$^2$

\item a.) $a=4$ cm, stranski rob $=6\sqrt{2}$ cm, $v=8$ cm \quad
      b.) $\cos\alpha=\dfrac{1}{3}$, $\alpha=70^{\circ}31'44''$ \quad
      c.) $V=\dfrac{128}{3}-\dfrac{1}{12}=\dfrac{511}{12}\approx 42{,}58$ cm$^3$

\item a.) $P=\dfrac{3}{4}\cdot 4\pi r^2+2\cdot\dfrac{\pi r^2}{2}=100\pi\approx 314{,}16$ dm$^2$ \quad
      b.) $V=125\pi\approx 392{,}70$ dm$^3$, $m=375\pi\approx 1178{,}1$ kg \quad
      c.) $a=\sqrt[3]{125\pi}\approx 7{,}323$ dm, $d=a\sqrt{3}\approx 12{,}68$ dm

\item a.) $v=10$ cm \quad
      b.) $S_0=1320$ cm$^2$, $V=13200$ cm$^3$ \quad
      c.) telesna diagonala $=\sqrt{3464}\approx 58{,}86$ cm, kot $=9^{\circ}46'57''$

\item a.) $v=4$ dm, $V=128\pi\approx 402{,}12$ dm$^3$, $P=120\pi\approx 376{,}99$ dm$^2$ \quad
      b.) $3\pi r^2=120\pi$, $r=2\sqrt{10}\approx 6{,}32$ dm

\item a.) $P=896$ cm$^2$, $V=1568$ cm$^3$ \quad
      b.) $\alpha=67^{\circ}35'6''$ \quad
      c.) $\omega=73^{\circ}44'23''$ \quad
      d.) $1372:196=7:1$

\item a.) $R=\dfrac{\sqrt{676}}{2}=13$ cm \quad
      b.) največja ploskev meri $24\times 8$; $2\pi r=24$, $r=\dfrac{12}{\pi}\approx 3{,}82$ cm

\item Iz $rv=12$ in $r^2+v^2=25$ dobimo dve rešitvi: \\
      $r=3,\ v=4$: $P=24\pi\approx 75{,}40$ dm$^2$, $V=12\pi\approx 37{,}70$ dm$^3$,
      središčni kot $216^{\circ}$; \\
      $r=4,\ v=3$: $P=36\pi\approx 113{,}10$ dm$^2$, $V=16\pi\approx 50{,}27$ dm$^3$,
      središčni kot $288^{\circ}$

\item $r=5$ cm, $v=10$ cm \quad
      a.) $V=\dfrac{500\pi}{3}\approx 523{,}60$ cm$^3$ \quad
      b.) $\dfrac{100\pi}{150\pi}=66{,}\overline{6}\ \%$

\item $R\approx 7{,}839$ dm, $S_0=72\left(1+\sqrt{2}\right)\approx 173{,}823$ dm$^2$,
      $V\approx 1362{,}667$ dm$^3$, $P\approx 723{,}937$ dm$^2$

\item a.) $v=5$ cm \quad
      b.) plašč $=130+120\sqrt{2}\approx 299{,}71$ cm$^2$ \quad
      c.) stranski rob $=\sqrt{194}\approx 13{,}93$ cm, $\alpha=21^{\circ}2'15''$

\item Stranica romba meri 13 cm, višina romba $\dfrac{120}{13}$ cm:
      $P=480\pi\approx 1507{,}96$ cm$^2$, \ $V=\dfrac{14400\pi}{13}\approx 3479{,}92$ cm$^3$

\item $S_0=36$ m$^2$, polmer včrtanega kroga $\varrho=2$ m \quad
      a.) $\dfrac{\pi\varrho^2}{S_0}=\dfrac{\pi}{9}\approx 34{,}91\ \%$ \quad
      b.) $P=144$ m$^2$

\item a.) $v=3$ cm, $V=18$ cm$^3$ \quad
      b.) $P=\dfrac{45}{2}+\dfrac{27\sqrt{3}}{2}\approx 45{,}88$ cm$^2$

\item Iz $\dfrac{1}{3}\pi r^2v=\dfrac{2}{3}\pi r^3$ sledi $v=2r$ in $s=r\sqrt{5}$. \quad
      a.) $r:s=1:\sqrt{5}$ \quad
      b.) $2\pi r^2:\pi rs=2:\sqrt{5}$

\item a.) $V_{valja}=1{,}728\pi^2$ dm$^3$, dvig $=\dfrac{1{,}728\pi^2}{32\pi^2}=0{,}054$ dm $=5{,}4$ mm \quad
      b.) $d=\sqrt{(8\pi)^2+(4\pi)^2+\pi^2}=9\pi\approx 28{,}27$ dm
\end{enumerate}
\normalsize

\end{document}
